\chapter{结论与展望}
\label{ch:conclusion}

\section{本文工作总结}
\label{sec:5_1}

本文面向复杂城市环境下手机直连卫星场景中的非合作终端定位识别需求，围绕低成本无人机平台条件下的非合作式感知与定位问题，系统开展了单目标自适应融合定位与多目标/协同鲁棒定位两方面研究。本文研究并不替代协议层身份认证或执法判定，而是面向异常接入、异常辐射或疑似违规使用线索触发后的终端定位排查环节。针对遮挡、多径、非视距传播、目标数量不确定以及无人机协同网络不稳定等复杂因素对定位性能带来的制约，论文从传播场建模、观测构造、智能决策到实验验证逐步展开，形成了由单目标连续定位向多目标并发定位、再向残缺观测条件下协同鲁棒定位递进的研究思路。

第一，针对复杂城市环境中 RSSI、AoA 和 PDR 等多源观测可靠性随遮挡、多径和目标运动而动态变化的问题，本文提出了基于物理感知深度强化学习的单目标自适应融合定位方法。该方法将连续定位任务建模为序列决策过程，在状态空间中引入 RMS 时延扩展、RSSI 方差、几何一致性特征以及历史融合权重等物理先验信息，并通过增量式动作空间和联合奖励设计实现融合权重的自适应调整。实验结果表明，所提方法能够有效刻画环境变化对观测质量的影响，提升系统对复杂传播环境的适应能力，相较于基线方法获得了更高的定位精度和更稳定的轨迹输出。

第二，针对多目标数量不确定、空间采样稀疏以及多径伪响应导致传统逐峰判定方法易产生虚警的问题，本文构建了基于空间信号表征的无线电地图与区域级判别方法。该方法将无人机扫描得到的离散 RSSI 观测映射为具有空间邻域结构的二维信号表征，从而为多目标区域识别、真实目标与伪响应区分以及后续位置估计提供统一输入基础。实验结果表明，基于无线电地图的空间信号表征有助于缓解稀疏采样和多径伪峰带来的局部多峰混淆，提高真实目标区域与虚假响应区域的区分能力。

第三，针对实际协同网络中无人机节点掉点、失效和链路中断造成的残缺观测问题，本文构建了基于知识蒸馏的教师—学生 GAT 鲁棒定位框架，以完整网络观测指导残缺网络学习，实现了对缺失观测信息的有效补偿。实验结果表明，教师—学生知识蒸馏机制能够有效减缓节点掉点带来的性能退化，提升多无人机残缺观测条件下的栅格定位精度和协同定位鲁棒性。

综上，复杂城市环境下的非合作终端定位识别问题不能简单视为静态参数估计问题，而应被看作一个具有显著环境动态性、观测异质性和协同网络不确定性的感知决策问题。本文围绕手机直连卫星场景下的非合作终端定位识别问题，构建了由单目标自适应融合定位到多目标区域级判别、再到掉点条件下协同鲁棒定位的渐进式技术路线。研究结果说明，将传播机理建模、空间信号表征与鲁棒学习机制结合，有助于提升低成本无人机平台在复杂城市环境中的定位精度、环境适应能力和工程鲁棒性，也为后续面向真实场景部署的非合作终端感知与定位研究提供了可借鉴的技术基础。

\section{未来研究展望}
\label{sec:5_2}

尽管本文围绕复杂城市环境下手机直连卫星非合作终端的定位识别问题开展了较为系统的研究，并取得了一定结果，但当前工作仍主要建立在仿真环境和相对理想化的任务设定基础上，在真实场景验证深度、多目标区域判别精细化能力以及动态协同网络复杂性刻画等方面仍存在进一步完善空间。结合全文两条研究主线，后续工作可在单目标自适应融合定位、多目标/协同鲁棒定位以及真实部署验证三个层面继续深化。

第一，本文当前研究仍主要建立在仿真环境和相对理想化的任务设定基础上，后续需进一步加强真实场景验证与工程化评估。复杂城市环境中的传播条件、目标行为模式以及无人机平台运行状态均具有较强不确定性，仿真模型虽能重现主要传播机理，但仍难完全覆盖真实部署中的复杂因素。后续可结合实地飞行实验、真实街区测量数据和多源传感器记录，对本文方法在真实应用条件下的泛化能力、鲁棒性与实时性进行更充分验证。

第二，围绕单目标自适应融合定位这一主线，后续可进一步向更高维、更复杂约束的定位任务拓展。例如，将当前二维平面条件下的融合方法扩展到三维空间定位框架，并进一步考虑高度信息、俯仰角观测以及复杂空间几何约束，以提升方法对实际非合作终端定位任务的适用性。

第三，围绕多目标/协同鲁棒定位这一主线，后续仍需进一步增强复杂场景下的区域判别能力和动态协同能力。一方面，可继续研究多目标数量估计、真实目标与多径鬼影的细粒度区分、稀疏采样条件下更高保真的无线电地图重构，以及区域级判别与位置解算的一体化设计；另一方面，可面向更一般的动态拓扑场景研究掉点与恢复交替、链路质量波动和部分节点观测失真条件下的在线协同定位，并进一步探索中间特征蒸馏、关系蒸馏和拓扑蒸馏等更丰富的知识迁移机制。

第四，本文所提出的方法不仅面向手机直连卫星场景下的非合作终端定位识别问题，也可进一步拓展至更广泛的无线感知与协同定位任务，如应急通信保障、复杂环境下的无线电源定位以及低空协同感知网络管理等方向，从而进一步提升相关研究的工程应用价值。

综上所述，围绕低成本无人机平台条件下非合作终端定位识别问题，后续研究仍需沿着单目标自适应融合定位与多目标/协同鲁棒定位两条主线持续推进，并进一步强化真实场景验证与工程化部署研究。随着无人机平台能力提升、无线感知技术发展以及智能学习方法不断演进，该方向在复杂环境目标感知、协同定位与智能决策等方面仍具有较大的研究空间和应用潜力。
